\begin{claim}{C.35}{Proved}{An obstruction to extending the fixed median prescription}
There is a continuous $\R^2$-action on the plane and a based
ambient-null loop in its centralizer such that the C.34 median
homeomorphisms of the closed nonopen locus do not extend
equivariantly over an adjacent open orbit at one parameter.
The closed restricted operator image on that open orbit is
the free proper translation group $\R^2$.
The original loop nevertheless has an explicit equivariant
null-homotopy on the entire plane.
Thus this example refutes extension of the fixed C.34 median
prescription; it is not a failure of the level-$0$ comparison.
\end{claim}
\begin{proof}
Use the closed simplex
\[
 \Delta=\{(x,y,z)\in\R_{\geq0}^3:x+y+z=1\}.
\]
For $(a,b)\in\R^2$ define
\[
 T_{a,b}(x,y,z)
   =\frac{(e^a x,e^b y,z)}{e^a x+e^b y+z}.
                                                        \tag{C.35.1}
\]
The denominator is positive on $\Delta$.
Direct substitution gives
$T_{a,b}T_{c,d}=T_{a+c,b+d}$ and
$T_{a,b}^{-1}=T_{-a,-b}$, so (C.35.1) is a continuous
action by homeomorphisms, fixing the three vertices and
preserving each open edge.
On the interior the coordinate homeomorphism
\[
 (x,y,z)\longmapsto(u,v)=\bigl(\log(x/z),\log(y/z)\bigr),
 \qquad
 (u,v)\longmapsto
       \frac{(e^u,e^v,1)}{1+e^u+e^v}
                                                        \tag{C.35.2}
\]
turns this action into $(u,v)\mapsto(u+a,v+b)$.
Its restricted image is therefore the closed free proper
translation group on $\R^2$.

Double $\Delta$ along its entire boundary. Since $\Delta$
is a closed disc, this double is a sphere. Apply (C.35.1)
on both copies; the formulas agree on their identified
boundary and hence give a continuous sphere action.
Remove the fixed vertex $(0,0,1)$. The remaining surface
$M$ is a plane, and the action restricts to it.
It is orientation preserving: each map preserves the
orientation on either triangle interior, and preserves
the two sheets of the oriented double.

Write $Q=\{T_{a,b}:(a,b)\in\R^2\}$ for this operator
image. Its orbits in $M$ are the two triangle interiors,
the three open edges, and the remaining two vertices.
Use the following coordinates on the edges:
\[
\begin{array}{c|c|c}
 \text{edge}&\text{coordinate in }\R&
                      \text{translation by }T_{a,b}\\ \hline
 y=0&\log(x/z)&a\\
 x=0&\log(y/z)&b\\
 z=0&\log(x/y)&a-b .
\end{array}                                               \tag{C.35.3}
\]
Thus the union $Y$ of the three edges and the two
vertices is exactly the closed nonopen orbit locus.

Let $w(t)=(a(t),b(t))$ be the piecewise linear based
loop with successive values
\[
       (0,0),\quad(1,1),\quad(1,0),\quad(0,0)
           \quad\text{at }t=0,\tfrac13,\tfrac23,1,
                                                        \tag{C.35.4}
\]
and put $f_t=T_{a(t),b(t)}$.
All its values commute with $Q$, and
\[
                  F_{\sigma,t}
                  =T_{(1-\sigma)a(t),(1-\sigma)b(t)}
                                                        \tag{C.35.5}
\]
is a based equivariant null-homotopy on the entire plane.
Continuity, including that of the inverse family, follows
from (C.35.1), which is already continuous on the compact
double and preserves the removed point.
The subgroup generated by the loop is exactly $Q$:
its first segment contains all diagonal translations
$(r,r)$ for $0\leq r\leq1$, and the differences of
points of the second segment contain all $(0,r)$
for $-1\leq r\leq1$. Integer sums and inverses give
both full coordinate axes and hence $\R^2$.
Consequently C.34 uses precisely the orbit locus $Y$
just described.

At $s=2/3$ the prefix in the C.27 formula ends at
$r=1-s=1/3$. By (C.35.3)--(C.35.4), the three
prefix translation intervals are, in their displayed
order,
\[
                         [0,1],\quad[0,1],\quad\{0\}.
                                                        \tag{C.35.6}
\]
At $t=2/3$, the current edge translations are
$(1,0,1)$. Taking the prescribed medians with
the prefix endpoints gives
\[
       \left.Q^{\mathrm{med}}_{2/3,\,2/3}\right|_Y:
           \quad\text{edge translations }(1,0,0).
                                                        \tag{C.35.7}
\]
The vertices stay fixed. These edge maps fit
continuously at the vertices, as also guaranteed
by C.34.

Every $Q$-equivariant homeomorphism from one triangle
interior to itself is a translation in (C.35.2).
Indeed, commutation with all translations gives
$h(u,v)=(u,v)+h(0,0)$. If it maps one interior to
the other, the same equation holds in the identical
coordinates on the target sheet. The sheet exchange
fixes the common boundary, so in either case any
continuous boundary extension has edge translation
parameters
\[
  R:\R^2\longrightarrow\R^3,\qquad
                    R(A,B)=(A,B,A-B).                  \tag{C.35.8}
\]
This is the exact restriction morphism.
Its image is the kernel of
\[
       \ell:\R^3\longrightarrow\R,\qquad
                         \ell(r_1,r_2,r_3)=r_1-r_2-r_3.
                                                        \tag{C.35.9}
\]
But $\ell(1,0,0)=1$, so (C.35.7) does not lie in
that image. It cannot extend equivariantly even
over a single adjacent interior. This proves
all assertions, while (C.35.5) gives the explicit
global contraction of the original loop.
\end{proof}

\paragraph{Executed endpoint certificate and scope.}
The standard-library rational calculation
\texttt{code/part\_b/c35\_median\_extension\_certificate.py}
computes the prefix extrema, the three median
values, the restriction matrix in (C.35.8),
and the nonzero obstruction (C.35.9).
Its JSON output records $(1,0,0)$ and residual $1$.
This certificate verifies the finite endpoint
calculation; the preceding simplex, doubling,
and restriction arguments prove the topological
statement.

C.34 is a homeomorphism construction on the whole
nonopen locus, and C.33 contracts the open-orbit
restrictions. C.35 proves that a global argument
must permit changing the nonopen-locus contraction
to satisfy the actual boundary restriction
morphisms of the open components. The remaining
problem is therefore the existence of compatible
equivariant contractions, not extension of the
particular median family without modification.
C.36 and C.37 meet the actual boundary equations on
the closed union $X_2$ of open orbits with
two-dimensional proper operator image, including
infinite accumulation. Extending this family
compatibly to the exterior of $X_2$ and the other
proper operator cases remains required.
The general C.12 comparison remains unresolved.
